% GENERATED FILE. Edit canonical lessons, then run npm run generate:books.
% canonical-source-hash: fnv1a64:2759b718848d6bae
% canonical-lessons: LA-C340-arma, LA-C340-eques, LA-C340-pedes, LA-C340-lorica, LA-C340-clipeus

\chapter{Arma}
\label{ch:la-arma}
\hlchaptermodality{340}

\begin{chapteropening}
\textbf{By the end of this chapter:} \emph{I can name arms and kinds of soldier: weapons, a horseman, a foot soldier, the breastplate and the round shield.}

Complete the last lesson of chapter 340: I can name arms and kinds of soldier: weapons, a horseman, a foot soldier, the breastplate and the round shield.
\end{chapteropening}

\section[arma, armōrum]{arma, armōrum — weapons, arms}

\label{lesson:LA-C340-arma}



\begin{quote}
Before the new one: say the Latin for a raid, then the Latin for to fight it out.
\end{quote}

\subsection*{You'll want to know: arma, armōrum}
\textbf{arma, armōrum} — \textquotedblleft{}weapons, arms\textquotedblright{}.

\begin{quote}
Mīlitēs arma capere iussī sunt.
\end{quote}

\begin{quote}
\emph{The soldiers were ordered to take up arms.}
\end{quote}

Neuter, plural only. Mostly armour and shields, the things fitted to the body. A famous epic opens with \emph{arma virumque canō}.

\begin{cousinweb}
\textbf{arma} is inherited from a Proto-Indo-European root meaning \textquotedblleft{}to fit together\textquotedblright{} (EtymDB): fitted gear. English \emph{arms} and \emph{army} come from it through French.
\end{cousinweb}

\begin{grammarlens}[title={What you've built}]
One more word for telling, explaining and describing at B1.
\end{grammarlens}

\subsection*{Your turn}
\begin{itemize}
\raggedright
  \item \emph{Say it:} \emph{arma, armōrum}
  \item \emph{Say it:} \emph{arma, armōrum}, once more
  \item \emph{Say it:} a sentence of your own with \emph{arma, armōrum}
  \item \emph{Recall:} say the Latin for a raid, then the Latin for to fight it out, then say \emph{arma, armōrum} again
\end{itemize}

\begin{tcolorbox}[breakable,colback=teal!4,colframe=teal!35!black,title={Before you move on}]
What does arma, armōrum mean? (\textquotedblleft{}Weapons, arms\textquotedblright{}.) Say it once more.
\end{tcolorbox}
\section[eques, equitis]{eques, equitis — a horseman, a cavalryman; a knight}

\label{lesson:LA-C340-eques}



\begin{quote}
Before the new one: say the Latin for to fight it out, then the Latin for weapons.
\end{quote}

\subsection*{You'll want to know: eques, equitis}
\textbf{eques, equitis} — \textquotedblleft{}a horseman, a cavalryman; a knight\textquotedblright{}.

\begin{quote}
Equitēs hostium ex silvā ērūpērunt.
\end{quote}

\begin{quote}
\emph{The enemy's cavalry burst out of the forest.}
\end{quote}

Masculine, 3rd declension, from \emph{equus} (taught). The \emph{equitēs} were also a rich class below the senators. English \emph{equestrian}.

\begin{cousinweb}
\textbf{eques} is built on \emph{equus}, \textquotedblleft{}horse\textquotedblright{} (EtymDB). English \emph{equestrian} comes from its family.
\end{cousinweb}

\begin{grammarlens}[title={What you've built}]
One more word for telling, explaining and describing at B1.
\end{grammarlens}

\subsection*{Your turn}
\begin{itemize}
\raggedright
  \item \emph{Say it:} \emph{eques, equitis}
  \item \emph{Say it:} \emph{eques, equitis}, once more
  \item \emph{Say it:} a sentence of your own with \emph{eques, equitis}
  \item \emph{Recall:} say the Latin for to fight it out, then the Latin for weapons, then say \emph{eques, equitis} again
\end{itemize}

\begin{tcolorbox}[breakable,colback=teal!4,colframe=teal!35!black,title={Before you move on}]
What does eques, equitis mean? (\textquotedblleft{}A horseman, a cavalryman; a knight\textquotedblright{}.) Say it once more.
\end{tcolorbox}
\section[pedes, peditis]{pedes, peditis — a foot soldier}

\label{lesson:LA-C340-pedes}



\begin{quote}
Before the new one: say the Latin for weapons, then the Latin for a horseman.
\end{quote}

\subsection*{You'll want to know: pedes, peditis}
\textbf{pedes, peditis} — \textquotedblleft{}a foot soldier\textquotedblright{}.

\begin{quote}
Peditēs lentē post equitēs ībant.
\end{quote}

\begin{quote}
\emph{The foot soldiers were marching slowly behind the cavalry.}
\end{quote}

Masculine, 3rd declension, from \emph{pēs} (taught). Mind the forms: \emph{pedes} 'foot soldier', \emph{pedēs} 'feet'.

\begin{cousinweb}
\textbf{pedes} is built on \emph{pēs}, \textquotedblleft{}foot\textquotedblright{} (Lewis \& Short): one who goes on foot. English \emph{pedestrian} comes from its family.
\end{cousinweb}

\begin{grammarlens}[title={What you've built}]
One more word for telling, explaining and describing at B1.
\end{grammarlens}

\subsection*{Your turn}
\begin{itemize}
\raggedright
  \item \emph{Say it:} \emph{pedes, peditis}
  \item \emph{Say it:} \emph{pedes, peditis}, once more
  \item \emph{Say it:} a sentence of your own with \emph{pedes, peditis}
  \item \emph{Recall:} say the Latin for weapons, then the Latin for a horseman, then say \emph{pedes, peditis} again
\end{itemize}

\begin{tcolorbox}[breakable,colback=teal!4,colframe=teal!35!black,title={Before you move on}]
What does pedes, peditis mean? (\textquotedblleft{}A foot soldier\textquotedblright{}.) Say it once more.
\end{tcolorbox}
\section[lōrīca, lōrīcae]{lōrīca, lōrīcae — a breastplate, a cuirass}

\label{lesson:LA-C340-lorica}



\begin{quote}
Before the new one: say the Latin for a horseman, then the Latin for a foot soldier.
\end{quote}

\subsection*{You'll want to know: lōrīca, lōrīcae}
\textbf{lōrīca, lōrīcae} — \textquotedblleft{}a breastplate, a cuirass\textquotedblright{}.

\begin{quote}
Lōrīca mīlitem ā sagittā servāvit.
\end{quote}

\begin{quote}
\emph{His breastplate saved the soldier from the arrow.}
\end{quote}

Feminine. The first ones were made of leather straps.

\begin{cousinweb}
\textbf{lōrīca} is built on \emph{lōrum}, \textquotedblleft{}a leather strap\textquotedblright{} (Lewis \& Short): a corselet of leather thongs.
\end{cousinweb}

\begin{grammarlens}[title={What you've built}]
One more word for telling, explaining and describing at B1.
\end{grammarlens}

\subsection*{Your turn}
\begin{itemize}
\raggedright
  \item \emph{Say it:} \emph{lōrīca, lōrīcae}
  \item \emph{Say it:} \emph{lōrīca, lōrīcae}, once more
  \item \emph{Say it:} a sentence of your own with \emph{lōrīca, lōrīcae}
  \item \emph{Recall:} say the Latin for a horseman, then the Latin for a foot soldier, then say \emph{lōrīca, lōrīcae} again
\end{itemize}

\begin{tcolorbox}[breakable,colback=teal!4,colframe=teal!35!black,title={Before you move on}]
What does lōrīca, lōrīcae mean? (\textquotedblleft{}A breastplate, a cuirass\textquotedblright{}.) Say it once more.
\end{tcolorbox}
\section[clipeus, clipeī]{clipeus, clipeī — a round shield}

\label{lesson:LA-C340-clipeus}



\begin{quote}
Before the new one: say the Latin for a foot soldier, then the Latin for a breastplate.
\end{quote}

\subsection*{You'll want to know: clipeus, clipeī}
\textbf{clipeus, clipeī} — \textquotedblleft{}a round shield\textquotedblright{}.

\begin{quote}
Clipeus aureus in templō pendēbat.
\end{quote}

\begin{quote}
\emph{A golden round shield hung in the temple.}
\end{quote}

Masculine (also neuter \emph{clipeum}). A round bronze shield; the long shield of the legions is the \emph{scūtum} (taught).

\begin{cousinweb}
\textbf{clipeus} is linked by Lewis \& Short with \emph{clepō}, \textquotedblleft{}to steal, to hide\textquotedblright{}: the cover. The link is not certain.
\end{cousinweb}

\begin{grammarlens}[title={What you've built}]
One more word for telling, explaining and describing at B1.
\end{grammarlens}

\subsection*{Your turn}
\begin{itemize}
\raggedright
  \item \emph{Say it:} \emph{clipeus, clipeī}
  \item \emph{Say it:} \emph{clipeus, clipeī}, once more
  \item \emph{Say it:} a sentence of your own with \emph{clipeus, clipeī}
  \item \emph{Recall:} say the Latin for a foot soldier, then the Latin for a breastplate, then say \emph{clipeus, clipeī} again
\end{itemize}

\begin{tcolorbox}[breakable,colback=teal!4,colframe=teal!35!black,title={Before you move on}]
What does clipeus, clipeī mean? (\textquotedblleft{}A round shield\textquotedblright{}.) Say it once more.
\end{tcolorbox}
